% !Mode:: "TeX:UTF-8"

\cabstract{
管道泄漏会影响能源输送、设备安全和环境稳定，传统检测方法往往依赖人工特征和阈值设定，在复杂噪声与长时序信号条件下泛化能力有限。针对上述问题，本文围绕双通道传感信号开展管道泄漏识别与定位研究。首先，依据泄漏声信号传播机理和左右通道响应差异，构建从原始 CSV 数据到固定窗口训练样本的处理流程，并采用文件级划分避免相邻片段信息泄漏。其次，建立两阶段学习框架：Stage2 使用单通道分段信号完成正常状态与多类泄漏形态识别，Stage1 使用左右双通道信号对进行泄漏距离回归。模型方面，本文以一维卷积网络为基础结构，引入 Conformer 模型结构对比，并设置传统机器学习基线和通道消融实验分析不同方法的适用边界。实验结果表明，一维卷积模型在 Stage2 测试集上取得片段级 Accuracy 0.9415、Macro-F1 0.9180，文件级 Accuracy 0.9783、Macro-F1 0.9836；Conformer 进一步将片段级 Accuracy 提升至 0.9900、Macro-F1 提升至 0.9883。Stage1 回归任务中，一维卷积模型取得片段级 MAE 1.90 m、RMSE 2.43 m，验证了双通道原始波形中包含可学习的距离信息，但传统基线和结构对比结果也表明，距离回归仍受样本规模和距离档位分布制约。本文工作验证了两阶段框架在泄漏识别与定位分析中的可行性。
}

\ckeywords{管道泄漏检测，时序信号分析，一维卷积神经网络，Conformer 模型，距离回归}

\eabstract{
Pipeline leakage may threaten energy transportation, equipment safety, and environmental stability. Traditional detection methods usually rely on handcrafted features and empirical thresholds, which limits their robustness under complex noise and long time-series conditions. This thesis studies a two-stage leak detection and localization method based on dual-channel sensor signals. First, according to the propagation mechanism of leakage acoustic signals and the response difference between left and right sensing channels, a data processing pipeline is built to convert raw CSV records into fixed-window training samples, and file-level splitting is adopted to avoid information leakage between adjacent segments. Second, two learning tasks are constructed: Stage2 uses segmented single-channel signals to recognize normal states and multiple leak patterns, while Stage1 uses paired dual-channel signals to estimate leak distance. A one-dimensional convolutional neural network is used as the basic model, and Conformer models, traditional machine-learning baselines, and channel ablation experiments are introduced for comparison. Experimental results show that the 1D CNN achieves a segment-level Accuracy of 0.9415 and Macro-F1 of 0.9180 on Stage2, with file-level Accuracy and Macro-F1 reaching 0.9783 and 0.9836. The Conformer model further improves the segment-level Accuracy to 0.9900 and Macro-F1 to 0.9883. On Stage1, the 1D CNN obtains a segment-level MAE of 1.90 m and RMSE of 2.43 m, indicating that distance-related information can be learned from raw dual-channel waveforms. However, the baseline and structure comparison results also show that distance regression is still limited by sample size and distance-bin distribution. The proposed framework verifies the feasibility of two-stage leak recognition and localization analysis.
}

\ekeywords{Leak Detection; Time-series; 1D CNN; Conformer; Regression}

\makeabstract
